% =============================================================================
% Chapter 8 - Screen Memory
% =============================================================================
\chapter{Screen Memory}

\epigraph{``The screen is a window into the machine's mind. What you write
there, the world will see.''}{\textit{--- An anonymous systems programmer}}

\section*{Introduction}

NovaVM's text layer is an 80$\times$50 display. Character RAM and Color RAM
are each 4000 bytes: one byte per visible cell. Character RAM holds the
character code; Color RAM holds the foreground colour nibble for the matching
cell.

These memories are no longer mapped directly into the 6502 address space.
The old \$AA00 and \$B1D0 windows are ordinary RAM now. All CPU access to
character and colour memory goes through the VDC-style VRAM port at
\$A0E0--\$A0E4. DMA and the blitter use the same memory-space IDs.

\begin{center}
\small
\begin{tabular}{@{}lll@{}}
\toprule
\textbf{Plane} & \textbf{Memory} & \textbf{Size} \\
\midrule
1 & Character RAM & 4000 bytes \\
2 & Color RAM & 4000 bytes \\
3 & Graphics bitmap & 64000 bytes \\
4 & Sprite shape RAM & 32768 bytes \\
6 & Tile data RAM & 32768 bytes \\
\bottomrule
\end{tabular}
\end{center}

\section{VRAM Port}
\label{sec:charram}

\memrange{A0E0}{A0E4}{VRAM Port}{R/W}{---}

\begin{center}
\small
\begin{tabular}{@{}lll@{}}
\toprule
\textbf{Address} & \textbf{Name} & \textbf{Purpose} \\
\midrule
\$A0E0 & VramPlane & Plane selector \\
\$A0E1 & VramAddrL & Low byte of byte offset \\
\$A0E2 & VramAddrH & High byte of byte offset \\
\$A0E3 & VramData & Data port \\
\$A0E4 & VramCtrl & Bit 0 = auto-increment after data access \\
\bottomrule
\end{tabular}
\end{center}

To access a byte, select the plane, write the 16-bit byte offset, configure
auto-increment if desired, then read or write \texttt{VramData}. Reads are
BRAM-backed and behave like a real VDC port: the first read posts the address,
and the following read returns the byte.

\section{Coordinate Math}
\label{sec:screen-math}

Both text planes use the same linear address formula:

\begin{center}
\texttt{offset = row * 80 + column}
\end{center}

\noindent where \texttt{row} ranges from 0 to 49 and \texttt{column} ranges
from 0 to 79.

\subsection{Worked Example}

To place the letter \texttt{A} (ASCII \$41) in colour 14 (light blue) at
column 35, row 10:

\begin{enumerate}
  \item Offset: 10 $\times$ 80 + 35 = 835 = \texttt{\$0343}
  \item Write \texttt{\$41} to plane 1 at offset \texttt{\$0343}
  \item Write \texttt{14} to plane 2 at offset \texttt{\$0343}
\end{enumerate}

\section{Assembly Example}

The following routine streams \texttt{HELLO} to the top-left corner using
auto-increment:

\begin{lstlisting}
VramPlane = $A0E0
VramAddrL = $A0E1
VramAddrH = $A0E2
VramData  = $A0E3
VramCtrl  = $A0E4

        lda #1              ; character plane
        sta VramPlane
        lda #0
        sta VramAddrL
        sta VramAddrH
        lda #1              ; auto-increment
        sta VramCtrl

        ldx #0
chars:  lda hello,x
        beq colors
        sta VramData
        inx
        bne chars

colors: lda #2              ; colour plane
        sta VramPlane
        lda #0
        sta VramAddrL
        sta VramAddrH
        lda #1
        sta VramCtrl

        ldx #0
color_loop:
        lda hello,x
        beq done
        lda #2              ; red
        sta VramData
        inx
        bne color_loop

done:   rts
hello:  .byte "HELLO",0
\end{lstlisting}

\section{Bulk Updates}

For small updates, the VRAM port is simple and predictable. For full-screen
or rectangular transfers, use DMA or blitter paths with the same plane IDs.
That keeps assembly code independent of internal BRAM layout and avoids
reintroducing wide CPU-visible memory windows.

\begin{warningbox}
\$AA00 and \$B1D0 are not screen memory on current NovaVM hardware. Code that
uses the old direct windows will read and write normal RAM and will not update
the display.
\end{warningbox}
